\reversemarginpar

\DeclareSourcemap{
    \maps[datatype=bibtex]{
        \map{
            \step[fieldsource=icon, match=\regexp{.}, final]
            \step[fieldset=usera, origfieldval]
        }
    }
}

% The icon field: a badge beside the entry's label. icon = {award} (or the older
% images/hm.png) draws the award rosette below; any other value is an image
% file. In the APR version, when the entry is highlighted (in the review
% period), the rosette is filled with the award colour, so an award on a
% paper or talk shows up as an award as well as in the entry's own colours.
% \cvawardbadgepaths{fill}{line width}: the rosette's paths, 20 x 22.5 units,
% for use inside a tikzpicture (the APR dots style draws it with its marker)
\newcommand\cvawardbadgepaths[2]{%
      \fill (6,13) -- (2.5,1) -- (5.5,3) -- (8,0) -- (10.5,10) -- cycle;
      \fill (14,13) -- (17.5,1) -- (14.5,3) -- (12,0) -- (9.5,10) -- cycle;
      \filldraw[fill=#1,line width=#2] (10,15.5) circle[radius=7];}
\newcommand\cvawardbadge[1]{% #1: fill colour of the medallion
    \begin{tikzpicture}[x=0.0625em,y=0.0625em,line join=round,
                        cv/decorative]% the award is in the entry's note
      \cvawardbadgepaths{#1}{0.14em}
    \end{tikzpicture}}
\newcommand\cvbadge[1]{%
    \aprsetentrycolour
    \ifboolexpr{test {\ifstrequal{#1}{award}} or test {\ifstrequal{#1}{images/hm.png}}}
      {\ifdefempty\aprentrycolour
         {\cvawardbadge{white}}%
         {\cvawardbadge{APRaward}}}%
      {\ifcv@tagged\includegraphics[height=1.5em,artifact]{#1}%
       \else\includegraphics[height=1.5em]{#1}\fi}}
% cvawardinline: set by the APR dots style, which draws the award badge as
% part of its marker; other badges stay in the margin
\newif\ifcvawardinline
% Otherwise the badge sits just left of the content column (\cvcontentindent,
% uwhealthcv.sty), where an entry's label ([J1]) starts, 6pt from it and
% centred on the first line, as the dots style places a rosette with no
% dots. An entry outside the content column keeps it in the margin.
\newlength\cv@badgeshift
\DeclareFieldFormat{usera}{%
    \ifboolexpr{bool {cvawardinline} and
        (test {\ifstrequal{#1}{award}} or test {\ifstrequal{#1}{images/hm.png}})}
      {}%
      {\setlength\cv@badgeshift{\dimexpr\@totalleftmargin
         \ifdefined\cvcontentindent-\cvcontentindent\fi+6pt\relax}%
       \ifdefined\cvcontentindent\ifdim\@totalleftmargin<\cvcontentindent
         \setlength\cv@badgeshift{-1pt}\fi\fi
       \ifdim\cv@badgeshift<\z@
         \marginpar{\hspace{1.75em}\raisebox{-0.5\height}{\cvbadge{#1}}}%
       \else
         % the rosette's medallion (69% up its height) on the line's middle,
         % ribbons hanging below; an image centred
         \ifboolexpr{test {\ifstrequal{#1}{award}} or test {\ifstrequal{#1}{images/hm.png}}}
           {\def\cv@badgeat{0.69}}{\def\cv@badgeat{0.5}}%
         \llap{\smash{\raisebox{\dimexpr0.2\baselineskip-\cv@badgeat\height\relax}{\cvbadge{#1}}}%
               \hspace{\cv@badgeshift}}%
       \fi}%
}

% The badge sits in the margin and takes no room in the line, so no space
% follows it: entries with and without a badge start at the same place.
\renewbibmacro*{begentry}{%
    \printfield{usera}%
}


% \setentrydetail{key}{text}: a line of small text printed under one
% bibliography entry, after its final full stop (e.g. a software entry's
% releases, commits and lines of code). Kept outside the .bib so a script can
% regenerate it without touching the entry. In the APR version the line is
% part of the entry, so it takes the entry's highlighting (\aprlinemark ends
% the first line's highlight at its text).
\newcommand\setentrydetail[2]{\csgdef{cv@detail@#1}{#2}}
\renewbibmacro*{finentry}{%
    \finentry
    \edef\cv@detailkey{\thefield{entrykey}}%
    \ifcsdef{cv@detail@\cv@detailkey}
      {\aprlinemark\newline{\small\csuse{cv@detail@\cv@detailkey}}}%
      {}}

% Redefine the bibliography driver for the "inproceedings" type
\DeclareBibliographyDriver{inproceedings}{%
  \usebibmacro{bibindex}%
  \usebibmacro{begentry}%
  \usebibmacro{author/editor}%
  \setunit{\labelnamepunct}\newblock
  \usebibmacro{title}%
  \newunit\newblock
  \printfield{booktitle}%
  \newunit\newblock
  \printfield{volume}%
  \printfield{number}%
  \printfield{series}%
  \newunit\newblock
  \printfield{pages}%
  \newunit\newblock
  \printfield{organization}% Print organization here
  \newunit\newblock
  \printfield{publisher}% Print publisher here
  \setunit{\addcomma\space}% Add comma before year
  \printfield{location}% Print location here
  \setunit{\addcomma\space}% Add comma before year
  \printfield{year}% Print year here
  \newunit\newblock
  \printfield{note}% Print note here
  \newunit\newblock
  \usebibmacro{doi+eprint+url}% Print DOI, eprint, and URL here
  \newunit\newblock
  \usebibmacro{finentry}%
}


\setlength{\marginparsep}{-10pt}

% A URL's break points don't stretch (biblatex's default is 0mu plus 3mu), so
% a justified line that ends in a URL stretches its word spaces instead of
% spreading the URL out ("https : / / example . org")
\biburlbigskip=0mu\relax









